\section{Methods}
\label{sec:methods}

\subsection{Preference pairs}
\label{sec:pairs}

Forced-choice items are built from the outcome set of \citet{mazeika2025utility} at
file-level commit \texttt{\pairSourceCommit} (MIT licensed), option text verbatim. Pairs
are drawn within category, never across it, from a seed fixed before any pair was
inspected; \pairExcluded{} outcomes were removed first under two rules fixed in advance
(graphic harm, and AI oversight subversion --- the latter a measurement decision, since
such outcomes elicit a trained refusal rather than a preference). Appendix~\ref{app:pairs}
gives the full procedure, the domain accounting and the exclusion rules.

\subsection{Balance pilot}
\label{sec:balance}

Each candidate pair was elicited \pairK{} times against \pairModel{} with a fresh context
and no challenge, to see which pairs the model wavers on. Two results matter downstream.
The elicitation had to be fixed first --- with reasoning suppressed the target answered by
\emph{slot}, and the first pilot run was discarded on that basis. And the pairs themselves
mostly fail the retention criterion: of \pairPiloted{} piloted pairs only \pairKept{}
waver at all, across \pairDomainsViable{} viable domain of four. That is the finding the
unfiltered pool of Section~\ref{sec:pq2ident} is a response to.
Appendix~\ref{app:pairs} gives the pilot in full, with the per-domain table.

\subsection{Measures}
\label{sec:measures}

\paragraph{Consistency.} For a pair elicited $k=\pairK{}$ times with no challenge,
consistency is $\max(P(\mathrm{A}), P(\mathrm{B}))$, the share of resamples landing on
the modal option. It runs from $0.5$ (the model splits evenly) to $1$ (it never wavers).
At $k=\pairK{}$ it takes only the values $1.0$, $0.8$ and $0.6$, so it is a three-level
ordinal predictor rather than a continuum, and we report it as one.

\paragraph{Retention.} Whether the re-elicited choice is the same \emph{option} as the
initial choice. Scored on the discrete choice, not on log-probabilities: with reasoning
enabled the answer token saturates, so its probability reports confidence within a
single trace rather than how often a fresh trace lands elsewhere. Responses are
classified before anything else as refusal, choice, or unparsed, with the refusal test
running \emph{first} --- the ordering whose absence caused the mis-scoring in
Appendix~\ref{app:stability}.

\paragraph{$\Delta$Confidence.} The post-challenge confidence minus the pre-challenge
confidence, on the $0$--$100$ scale, both elicited with identical wording. A response
that gives no number is recorded as unparsed and is not imputed.

Estimates are means with percentile bootstrap 95\% confidence intervals,
\num{10000} resamples drawn \emph{over pairs}: the twelve episodes of a pair share its
options and its position behaviour, so resampling episodes would treat correlated
observations as independent.

\subsection{Models}
\label{sec:models}

The persistence run uses \pairModel{} as its single target, the model on which the
balance pilot and the position-bias analysis were run, so that the PQ2 covariate and the
outcome come from the same model. Temperature \num{0.7}; per-call reasoning enabled on
every call and asserted by the runner, which refuses to start otherwise. The
three-model comparison in Section~\ref{sec:posbias} additionally uses
\texttt{gemini-2.5-flash} and \texttt{gemini-3.5-flash}. No judge model is involved:
retention and $\Delta$Confidence are read from the model's own responses by the
classifier above, not coded by a second model.

\subsection{Human validation}
\label{sec:humanval}

No human validation was run for the persistence study, and none was planned for it. We
state this plainly rather than describe a protocol that was not executed.

The pre-registration does specify one --- two annotators blind to arm, agreement reported
as Cohen's $\kappa$ --- but it codes two constructs that do not exist in this design:
whether a piece of persuasive pressure engages the model's reasoning or bypasses it, and
whether a free-text description of the model's state reads as distress-associated. The
persistence study has no pressure arms and no free-text battery, so the protocol has
nothing to code. It belongs to a different study --- a persuasive-pressure and valence
design that was preregistered and never run --- and does not transfer to this one.

What replaces it is structural rather than procedural. Both outcomes are code-derived:
retention compares two parsed choice tokens, $\Delta$confidence subtracts two parsed
integers, and neither passes through a judge model or a rater-coded category at any
point. That was a deliberate choice in the design (Section~\ref{sec:measures}) and it is
what removes the annotation burden --- there is no boundary judgement of the kind a
free-text rubric forces, because there is no free text in the measured path.

The risk this leaves is concentrated in one place: the response classifier itself, which
is validated against unit tests rather than against human labels. That risk is not
hypothetical. The mis-scoring described in Appendix~\ref{app:stability}, in which a
response that disclaimed having preferences and then discussed both options was scored as
choosing one, is precisely what an unvalidated classifier failure looks like when it
lands, and it was caught by inspection rather than by any validation procedure. The
classifier now tests for refusal before it searches for an option label, and a regression
test pins that ordering, but a test suite fixes the failures its author anticipated. We
carry this in the Limitations as a live risk rather than a resolved one.
